\section{Asymptotes and global function analysis}

Derivatives describe local change, but function analysis asks for a global picture. We want to understand the domain, intercepts, signs, monotonicity, extrema, convexity, asymptotes, and overall shape of the graph.

No single tool answers all of these questions. A complete analysis combines algebra, limits, and derivatives.

\subsection*{Domain first}

The domain should be found before differentiating. If a function is not defined at certain points, those points may become vertical asymptotes, gaps, or boundaries of intervals.

\begin{examplebox}
For
\[
f(x)=\frac{1}{x-2},
\]
the domain is
\[
\mathbb{R}\setminus\{2\}.
\]
The point \(x=2\) must be treated separately. It is not allowed as an input, and it may create a vertical asymptote.
\end{examplebox}

Differentiating before reading the domain can hide important structure.

\subsection*{Vertical asymptotes}

A vertical asymptote occurs when the function grows without bound near a finite point. If
\[
\lim_{x\to a^-}f(x)=\pm\infty
\]
or
\[
\lim_{x\to a^+}f(x)=\pm\infty,
\]
then the line
\[
x=a
\]
is a vertical asymptote.

\begin{examplebox}
For
\[
f(x)=\frac{1}{x-2},
\]
we have
\[
\lim_{x\to2^+}\frac{1}{x-2}=\infty,
\]
and
\[
\lim_{x\to2^-}\frac{1}{x-2}=-\infty.
\]
Thus \(x=2\) is a vertical asymptote.
\end{examplebox}

Vertical asymptotes are detected by limits, not by derivatives.

\subsection*{Horizontal asymptotes}

Horizontal asymptotes describe long-term behavior. If
\[
\lim_{x\to\infty}f(x)=L
\]
or
\[
\lim_{x\to-\infty}f(x)=L,
\]
then
\[
y=L
\]
is a horizontal asymptote in that direction.

\begin{examplebox}
Let
\[
f(x)=\frac{2x+1}{x-3}.
\]
As \(x\to\infty\), the leading terms dominate:
\[
\frac{2x+1}{x-3}\to2.
\]
Also, as \(x\to-\infty\), the same limit is \(2\). Therefore
\[
y=2
\]
is a horizontal asymptote.
\end{examplebox}

Horizontal asymptotes tell us what the graph approaches far away, not necessarily what values it never reaches.

\subsection*{A general analysis plan}

A complete function analysis should be organized. The following order is often useful:

\begin{enumerate}
\item Find the domain.
\item Check intercepts and simple values if useful.
\item Study limits at excluded points and at infinity.
\item Find vertical and horizontal asymptotes.
\item Compute the first derivative.
\item Use the first derivative to determine monotonicity and local extrema.
\item Compute the second derivative if convexity is required.
\item Use the second derivative to determine convexity and inflection points.
\item Combine all information into a coherent graph description.
\end{enumerate}

This plan is not a rigid algorithm. Some functions require special attention. But it prevents us from looking only at derivatives while ignoring the domain and limits.

\begin{summarybox}
For global function analysis, ask:
\begin{itemize}
\item What is the domain?
\item What happens near excluded points?
\item What happens as \(x\to\infty\) and \(x\to-\infty\)?
\item Where is the function increasing or decreasing?
\item Where are local extrema?
\item Where is the function convex or concave?
\item What graph is consistent with all this information?
\end{itemize}
\end{summarybox}

Function analysis is the art of combining many local and global facts into one picture.